\paragraph{Truncated moment certificates and Markov--Stieltjes extremals: uncertainty intervals under finite register length and endpoint realization}
Let \(\mu\) be a finite positive measure supported on \([1,\infty)\), and let its Stieltjes orbit be
\[
S(t):=\int_{[1,\infty)}\frac{1}{t+a}\,\dd\mu(a),\qquad t\ge 0.
\]
Since the spectral gap \(a\ge 1\), the function \(S\) is analytic in a neighborhood of \(t=0\) and induces the truncated moment chain (namely, the Taylor coefficients of \(S\) at \(t=0\))
\[
\mu_n:=\frac{(-1)^n}{n!}S^{(n)}(0)=\int_{[1,\infty)}a^{-(n+1)}\,\dd\mu(a),\qquad n\ge 0.
\]

\begin{proposition}[Markov--Stieltjes extremal principle: strongest feasible interval under truncated moment information and endpoint realization]\label{prop:xi-markov-stieltjes-extremal-principle}
Fix an integer \(N\ge 1\) and given a truncated moment vector \((\mu_0,\dots,\mu_{2N-1})\).
Let \(\mathfrak{M}_N\) be the set of all finite positive measures satisfying the above \(2N\) moment constraints:
\[
\mathfrak{M}_N:=\Bigl\{\nu\ \text{is a finite positive measure on \([1,\infty)\)}:\ \int a^{-(n+1)}\,\dd\nu(a)=\mu_n\ (0\le n\le 2N-1)\Bigr\}.
\]
Then for any fixed \(t>0\), the feasible value set
\[
\bigl\{S_\nu(t):=\int\frac{1}{t+a}\,\dd\nu(a)\ :\ \nu\in\mathfrak{M}_N\bigr\}
\]
is a closed interval \([S_N^{-}(t),S_N^{+}(t)]\).
Here the endpoint functions \(S_N^\pm\) are two Stieltjes-type rational functions uniquely determined by \((\mu_0,\dots,\mu_{2N-1})\), and can be taken as adjacent convergents of the same positive-coefficient Stieltjes continued fraction (even/odd extremal branches).
Furthermore, the extremal measures attaining the upper and lower endpoints must be atomic measures with at most \(N\) atoms.
\end{proposition}
\begin{proof}[Proof outline]
This is the standard extremal conclusion for the truncated Stieltjes moment problem (Markov theorem/endpoint case of Ne\-van\-lin\-na parametrization).
For fixed \(t>0\), the mapping \(\nu\mapsto S_\nu(t)\) is a continuous linear functional; by Nevanlinna parametrization, the set of all feasible values is precisely a closed interval, whose endpoints are realized by extreme solutions.
The extreme solutions of the truncated moment problem must be discrete measures with at most \(N\) atoms, and correspond to the two adjacent convergents (even/odd extremal branches) of the positive-coefficient continued fraction constructed by the Stieltjes algorithm from \((\mu_0,\dots,\mu_{2N-1})\). This completes the proof.
\end{proof}

\begin{proposition}[Minimum atomic realization of truncated moments: provable lower bound of defect degree and Gauss-type reconstruction]\label{prop:xi-truncated-moment-min-atomic-realization}
Under the setting of Proposition \ref{prop:xi-markov-stieltjes-extremal-principle}, denote the Hankel matrix block
\[
H_{N-1}:=\bigl(\mu_{i+j}\bigr)_{0\le i,j\le N-1}\succeq 0,
\qquad
\mathfrak{d}_N:=\mathrm{rank}(H_{N-1})\le N.
\]
Then there exist a set of nodes \(\{a_j^{(N)}\}_{j=1}^{\mathfrak{d}_N}\subset[1,\infty)\) and weights \(\{w_j^{(N)}\}_{j=1}^{\mathfrak{d}_N}\subset(0,\infty)\) such that
\[
\mu_n=\sum_{j=1}^{\mathfrak{d}_N} w_j^{(N)}\,\bigl(a_j^{(N)}\bigr)^{-(n+1)},
\qquad n=0,1,\dots,2N-1.
\]
Moreover, this minimum atomic realization can be constructed from the orthogonal polynomial system generated by the truncated moments: the nodes are the real zeros (lying in \((1,\infty)\)) of the corresponding orthogonal polynomial, and the weights are determined in closed form by the Christoffel numbers.
Therefore, the "defect degree" has a truncated version beyond the full-information limit \(\kappa\): given a finite certificate length \(2N\), the minimum realizable defect degree within the regime is at most \(N\), and can be constructively recovered from purely algebraic objects.
\end{proposition}
\begin{proof}[Proof outline]
From the truncated moments \((\mu_0,\dots,\mu_{2N-1})\) one can define a positive definite/semi-definite inner product on \(\mathrm{span}\{1,a^{-1},\dots,a^{-N}\}\), and obtain orthogonal polynomials and three-term recurrence; the spectrum of the corresponding truncated Jacobi matrix gives the nodes, while the first coordinate weights give the Christoffel numbers.
\(\mathfrak{d}_N=\mathrm{rank}(H_{N-1})\) is precisely the minimum atomic number lower bound for this finite-dimensional moment problem, and attainability is given by Gauss-type quadrature/spectral measure construction. This completes the proof.
\end{proof}

\begin{proposition}[Depth spectral Hankel determinant closed form: Vandermonde square law and algebraic closure of separation pathology]\label{prop:xi-depth-hankel-determinant-vandermonde-square}
Under the finite atomic regime, suppose
\[
\mu=\sum_{j=1}^{\kappa}w_j\,\delta_{a_j},
\qquad
1<a_1<a_2<\cdots<a_\kappa,\qquad w_j>0,
\]
and let the moment chain \(\mu_n=\int a^{-(n+1)}\,\dd\mu(a)\).
Denote \(b_j:=a_j^{-1}\in(0,1)\), \(\widetilde w_j:=w_j b_j>0\), then
\[
\mu_n=\sum_{j=1}^{\kappa}\widetilde w_j\,b_j^{n}\qquad(n\ge 0).
\]
Let the Hankel principal block \(\mathsf H_{\kappa-1}:=(\mu_{p+q})_{0\le p,q\le \kappa-1}\).
Then the Vandermonde factorization holds
\[
\boxed{\;
\mathsf H_{\kappa-1}=V\,\mathrm{diag}(\widetilde w_1,\dots,\widetilde w_\kappa)\,V^{\mathsf T},\qquad
V_{pj}:=b_j^{p}\;}
\]
and thus the determinant closed form
\begin{equation}\label{eq:xi-depth-hankel-det-vandermonde}
\boxed{\;
\det \mathsf H_{\kappa-1}
=\left(\prod_{j=1}^{\kappa}\widetilde w_j\right)
\left(\prod_{1\le j<k\le \kappa}(b_k-b_j)^2\right).\;}
\end{equation}
In particular, \(\det \mathsf H_{\kappa-1}=0\) if and only if there exists a spectral collision \(b_j=b_k\) or weight degeneration \(\widetilde w_j=0\).
Moreover, any local stability constant based on Hankel inversion (such as Prony/Pad\'e/Hankel null-space methods) cannot remain uniformly bounded at the \(b\)-collision limit: when \(\prod_{j<k}|b_k-b_j|\downarrow 0\), the error amplification factor diverges at least at a rate \(\bigl(\prod_{j<k}|b_k-b_j|^{-2}\bigr)\).
\end{proposition}

\begin{proof}
From
\(
\mu_{p+q}=\sum_j \widetilde w_j\,b_j^{p+q}
\)
we directly obtain the matrix factorization \(\mathsf H_{\kappa-1}=V\,\mathrm{diag}(\widetilde w)\,V^{\mathsf T}\).
Taking determinants and using \(\det(V)=\prod_{j<k}(b_k-b_j)\) yields \eqref{eq:xi-depth-hankel-det-vandermonde}.
The last statement is a standard condition number lower bound: \(\|\mathsf H_{\kappa-1}^{-1}\|\) is at least of the same scale as \((\det\mathsf H_{\kappa-1})^{-1}\) (up to a polynomial factor controlled by \(\|\mathsf H_{\kappa-1}\|\)), so collapse of the Vandermonde square factor must trigger inversion pathology explosion. This completes the proof.
\end{proof}

\begin{proposition}[Hankel determinant covariance law under scaling renormalization: scale factors can be explicitly extracted]\label{prop:xi-hankel-det-scaling-covariance}
Under the finite atomic setting of Proposition \ref{prop:xi-depth-hankel-determinant-vandermonde-square}, write \(a_j=1+\delta_j\) (\(\delta_j>0\)).
For any scaling parameter \(m>0\), define the scaled depth nodes
\[
\widetilde a_j:=1+m\delta_j,
\qquad
\widetilde b_j:=\widetilde a_j^{-1}=\frac{1}{1+m\delta_j},
\qquad
\widetilde{\widetilde w}_j:=w_j\widetilde b_j,
\]
and let \(\widetilde{\mathsf H}_{\kappa-1}\) be the Hankel principal block induced by the scaled moment chain
\(
\widetilde\mu_n:=\sum_{j=1}^{\kappa}w_j\,\widetilde a_j^{-(n+1)}
=\sum_{j=1}^{\kappa}\widetilde{\widetilde w}_j\,\widetilde b_j^{n}
\).
Then there exists an explicit covariance factor \(J_{\kappa-1}(m;\delta)>0\) such that
\[
\boxed{\;
\det \widetilde{\mathsf H}_{\kappa-1}
=\det \mathsf H_{\kappa-1}\cdot J_{\kappa-1}(m;\delta),\;}
\]
and one can take
\[
\boxed{\;
J_{\kappa-1}(m;\delta)
=m^{\kappa(\kappa-1)}
\prod_{j=1}^{\kappa}\Bigl(\frac{1+\delta_j}{1+m\delta_j}\Bigr)^{2\kappa-1}.\;}
\]
Therefore the scaling flow does not change the algebraic degeneration type of "collision/merging": when there exists \(\delta_j\to\delta_k\) (equivalent to \(b_j\to b_k\)), both \(\det \mathsf H_{\kappa-1}\) and \(\det\widetilde{\mathsf H}_{\kappa-1}\) exhibit second-order vanishing by the Vandermonde square law (Proposition \ref{prop:xi-depth-hankel-determinant-vandermonde-square}), while the scaling only introduces separable single-point scale factors through \(J_{\kappa-1}\).
\end{proposition}
\begin{proof}[Proof outline]
By \eqref{eq:xi-depth-hankel-det-vandermonde}, for scaled nodes we have
\[
\det \widetilde{\mathsf H}_{\kappa-1}
=\Bigl(\prod_{j=1}^{\kappa}w_j\widetilde b_j\Bigr)\Bigl(\prod_{1\le j<k\le\kappa}(\widetilde b_k-\widetilde b_j)^2\Bigr).
\]
Now
\(
\widetilde b_k-\widetilde b_j=\frac{1}{1+m\delta_k}-\frac{1}{1+m\delta_j}
=\frac{m(\delta_j-\delta_k)}{(1+m\delta_j)(1+m\delta_k)}
\).
Comparing this with \(b_k-b_j=\frac{\delta_j-\delta_k}{(1+\delta_j)(1+\delta_k)}\) and combining the products over all pairs \((j,k)\), we obtain the stated \(J_{\kappa-1}(m;\delta)\). This completes the proof.
\end{proof}

\begin{proposition}[Order preservation of truncated Pad\'e--Stieltjes approximation: poles are real negative, residues are positive, and yield monotone convergence across all scales]\label{prop:xi-pade-stieltjes-order-preserving}
If \(S\) is a Stieltjes function as above, then its diagonal Pad\'e approximant \([N/N]\) at \(t=0\) (matching the first \(2N\) Taylor coefficients) is still a Stieltjes-type rational function: all its poles lie in \((-\infty,-1)\) and are simple, and the residues in the corresponding partial fraction expansion are strictly positive.
Moreover, for any \(t>0\), the sequence \([N/N](t)\) splits into two monotone chains by parity in \(N\), squeezing the true value \(S(t)\) from above and below; and as \(N\to\infty\) the squeezing interval converges to the singleton \(\{S(t)\}\).
This conclusion elevates the auditability of "finite register certificates" to order preservation: the approximation process generates neither spurious poles nor sign-flipping spurious weights, thus excluding at the regime level any non-physical defect indices introduced by approximation.
\end{proposition}
\begin{proof}[Proof outline]
The fact that Pad\'e approximants of the Stieltjes class preserve positive weights and real negative pole structure is a direct consequence of continued fraction/orthogonal polynomial theory; the parity monotonicity corresponds to the alternating squeezing property of Stieltjes continued fraction convergents, and convergence comes from complete monotonicity and compactness control of the Stieltjes transform on \((0,\infty)\). This completes the proof.
\end{proof}

\begin{proposition}[Pure tail identifiability of the shallowest depth atom: moment ratio limit determines \texorpdfstring{$a_{\min}$}{a_min} and gives exponential convergence rate]\label{prop:xi-moment-ratio-identifies-amin}
Under the finite atomic regime, suppose
\(
\mu=\sum_{j=1}^{\kappa} w_j\,\delta_{a_j}
\)
with \(1<a_1:=a_{\min}<a_2\le\cdots\le a_{\kappa}\), \(w_j>0\).
Let the moment chain
\[
\mu_n:=\int a^{-(n+1)}\,\dd\mu(a)=\sum_{j=1}^{\kappa} w_j\,a_j^{-(n+1)}\qquad(n\ge 0).
\]
Then the limits exist
\[
\boxed{\;
\lim_{n\to\infty}\frac{\mu_n}{\mu_{n+1}}=a_1,
\qquad
\lim_{n\to\infty}\frac{\mu_{n+1}}{\mu_n}=\frac{1}{a_1}\; }.
\]
Moreover, the convergence is exponential: there exists a constant \(C>0\) such that
\[
\boxed{\;
\left|\frac{\mu_n}{\mu_{n+1}}-a_1\right|
\ \le\ C\left(\frac{a_1}{a_2}\right)^{n}\qquad(n\to\infty)\; }.
\]
Therefore the shallowest depth \(\delta_{\min}=a_1-1\), under the condition of not introducing any phase carrier and not performing full spectral inversion, can be uniquely recovered from the ratio limit of the single-variable moment chain tail; the spectral gap ratio \(a_1/a_2\) precisely controls its regime-level identifiability speed.
\end{proposition}
\begin{proof}
Write the moment chain as \(\mu_n=w_1a_1^{-(n+1)}(1+R_n)\), where
\(
R_n:=\sum_{j=2}^{\kappa}\frac{w_j}{w_1}\bigl(\frac{a_1}{a_j}\bigr)^{n+1}
\).
Then \(|R_n|\le C_0(\frac{a_1}{a_2})^{n}\) (for some constant \(C_0>0\)), and
\[
\frac{\mu_n}{\mu_{n+1}}
=a_1\cdot\frac{1+R_n}{1+R_{n+1}}
=a_1\left(1+\frac{R_n-R_{n+1}}{1+R_{n+1}}\right).
\]
Since \(R_n-R_{n+1}=O((a_1/a_2)^n)\) we obtain the stated exponential error bound. This completes the proof.
\end{proof}

\begin{proposition}[Tail rate limit law of Laplace partition: \texorpdfstring{$\delta_{\min}$}{delta_min} is the rate of the partition function]\label{prop:xi-partition-tail-rate-deltamin}
Under the finite atomic setting of the previous proposition, let \(\delta_j:=a_j-1\), and define the depth partition function
\[
Z_\delta(u):=\sum_{j=1}^{\kappa} w_j\,e^{-u\delta_j}\qquad(u\ge 0).
\]
Then the limit exists
\[
\boxed{\;
\lim_{u\to\infty}-\frac{1}{u}\log Z_\delta(u)=\delta_{\min}:=\min_j \delta_j\; }.
\]
Moreover, if \(\delta_2\) is the second smallest depth, then there is an exponential error bound: there exists a constant \(C>0\) such that
\[
\left|-\frac{1}{u}\log Z_\delta(u)-\delta_{\min}\right|
\ \le\ \frac{1}{u}\log\!\left(1+C\,e^{-u(\delta_2-\delta_{\min})}\right).
\]
This conclusion elevates the identifiability of "shallowest depth" to a thermodynamic limit law: \(\delta_{\min}\) is the rate function of the partition function, and the depth spectral gap \(\delta_2-\delta_{\min}\) determines the exponential scale of its convergence.
\end{proposition}
\begin{proof}
Write
\(
Z_\delta(u)=e^{-u\delta_{\min}}\bigl(w_1+\sum_{j\ge 2} w_j e^{-u(\delta_j-\delta_{\min})}\bigr)
\),
take logarithm and divide by \(u\) to obtain the limit.
The error bound is given by the exponential decay of the second term inside the bracket. This completes the proof.
\end{proof}

\begin{corollary}[Minimum observation scale of tail rate certificate: \texorpdfstring{$\delta_{\min}$}{delta_min}-driven indistinguishability lower bound]\label{cor:xi-partition-tail-min-scale-lower-bound}
Still under the finite defect setting of Proposition \ref{prop:xi-partition-tail-rate-deltamin}, suppose the regime only allows observation of the partition tail rate \(Z_\delta(u)\) (or equivalently observation of \(-\frac{1}{u}\log Z_\delta(u)\)) on a discrete grid
\(
u\in\{0,\Delta u,2\Delta u,\dots,M\Delta u\}
\),
and the observation accuracy requirement is to achieve relative error on the order of \(\varepsilon\in(0,1)\).
Then to distinguish in the worst case between "critical defect-free (\(\delta_{\min}=0\) in the closure sense)" and "existence of shallowest depth \(\delta_{\min}>0\)", the required maximum observation scale must satisfy
\[
\boxed{\;
M\Delta u\ \gtrsim\ \frac{1}{\delta_{\min}}\log\!\frac{1}{\varepsilon}\;}
\]
where the implicit constant depends only on a priori bounds on weight ratios (e.g., an upper bound on \(w_2/w_1\)).
This lower bound characterizes the invisible boundary layer within the regime: the smaller \(\delta_{\min}\) is, the larger \(u\) (longer "regime time") one must push the observation to in order to produce auditable separation.
\end{corollary}
\begin{proof}[Proof outline]
By \(Z_\delta(u)\ge w_1e^{-u\delta_{\min}}\) and \(Z_\delta(u)\le (w_1+\cdots+w_\kappa)e^{-u\delta_{\min}}\) we have: if \(\delta_{\min}=0\) (closure regime) then \(Z_\delta(u)\ge w_1\) does not decay with \(u\); whereas if \(\delta_{\min}>0\) then \(Z_\delta(u)\le \Phi\,e^{-u\delta_{\min}}\) decays exponentially (\(\Phi:=\sum_j w_j\)).
Therefore to distinguish the two under relative error \(\varepsilon\), one must enter the exponential tail regime \(e^{-u\delta_{\min}}\lesssim\varepsilon\), i.e., \(u\gtrsim \delta_{\min}^{-1}\log(1/\varepsilon)\).
If \(u\) can only reach \(M\Delta u\), we obtain the stated lower bound. This completes the proof.
\end{proof}

\endinput
